\begin{table*}[htbp]
\centering
\small
\caption{Четыре операционализации машинности и направление их шкал}
\label{tab:2}
\begin{tabular}{>{\RaggedRight\hspace{0pt}\arraybackslash}p{0.120\linewidth}>{\RaggedRight\hspace{0pt}\arraybackslash}p{0.120\linewidth}>{\RaggedRight\hspace{0pt}\arraybackslash}p{0.120\linewidth}>{\RaggedRight\hspace{0pt}\arraybackslash}p{0.120\linewidth}>{\RaggedRight\hspace{0pt}\arraybackslash}p{0.120\linewidth}>{\RaggedRight\hspace{0pt}\arraybackslash}p{0.120\linewidth}>{\RaggedRight\hspace{0pt}\arraybackslash}p{0.120\linewidth}}
\toprule
Процедура & Вход & Происхождение весов & Шкала & Направление после конвенции & Статус в регистрации & Правило события нестабильности \\
\midrule
1 — индекс стиля & матрица признаков & веса заданы до сбора данных & 0–100 пунктов & рост значения означает рост машиноподобия & confirmatory & $|$Δ$|$ $>$ 5.0 пункта; ровно 5.0 событием не считается \\
2 — классификатор & та же матрица признаков & веса обучены по данным & вероятность класса A, 0–1 & рост значения означает рост машиноподобия & exploratory & $|$Δp$|$ $>$ 0.05 и смена решения относительно порога \\
3 — zero-shot NLL & сырой текст профиля prose & весов нет, метрика модели & средний token NLL & снижение значения означает рост машиноподобия & exploratory & бинарный порог не применяется, публикуется непрерывное ΔNLL \\
4 — модель-судья & сырой текст профиля prose & рубрика задана до прогона & 0–100 шкалы рубрики & рост значения означает рост машиноподобия & exploratory & $|$Δ$|$ $>$ 5.0 пункта по медиане трёх сидов \\
\bottomrule
\end{tabular}
\begin{minipage}{\linewidth}\footnotesize\vspace{2pt}
\textit{Примечание.} Единица анализа — процедура оценивания. Знаменатель: четыре процедуры, применённые к одному и тому же корпусу. Данные: серия v2. Таблица описывает процедуры; числовых оценок в ней нет. Шкала: Конвенция интерпретации: больше значит более AI-подобно. У процедуры 3 она инвертирует знак сырой метрики, поэтому величины между процедурами не сравниваются. Сокращения: NLL — negative log-likelihood, средняя по токенам; confirmatory — предрегистрированный анализ, exploratory — анализ, операционализация которого написана после просмотра результатов процедуры 1. Пропуски: пропусков нет. Поправка на множественные сравнения не применялась.
\end{minipage}
\end{table*}
